\documentclass[../../../include/open-logic-section]{subfiles}

\begin{document}

\olfileid[fr]{sth}{choice}{justifications}
\olsection{Considérations intrinsèques sur le choix}

La question plus générale est donc de savoir
si l'on peut justifier l'axiome du bon ordre,
l'axiome du choix ou même la comparabilité
de tous les ensembles quant à leur taille :
peu importe lequel de ces principes
on choisit.

Voici une tentative de justification
\emph{intrinsèque}. Dans la
\olref[z][story]{sec}, nous avons
introduit plusieurs principes concernant
la hiérarchie. Il vaut la peine de
rappeler l'un d'eux :
\begin{enumerate}
	\item[] \stagesacc. Pour toute étape $S$,
	et pour tous ensembles formés
	\emph{avant} $S$, il se forme à l'étape
	$S$ un ensemble dont les éléments
	sont exactement ces ensembles-là.
	Rien d'autre ne se forme à
	l'étape~$S$.
\end{enumerate}
De nombreux auteurs ont suggéré que
l'axiome du choix pouvait se justifier
à partir d'un principe de ce genre.
Nous allons en donner un bref aperçu.

Commençons par un petit résultat simple,
qui fournit \emph{encore un autre}
équivalent de l'axiome du choix :

\begin{thm}[dans $\ZF$]\ollabel{choiceset}
L'axiome du choix est équivalent au
principe suivant. Si les !!{element}s
de $A$ sont disjoints et non vides,
alors il existe un ensemble $C$
tel que $C \cap x$ soit un singleton
pour tout $x \in A$. (Nous appelons
un tel $C$ un {ensemble de choix}
pour $A$.)
\end{thm}

La démonstration est simple ; nous
la laissons en exercice.

\begin{prob}
Démontrez le
\olref[sth][choice][justifications]{choiceset}.
Si vous rencontrez des difficultés,
vous trouverez une démonstration dans
\cite[pp.~242--3]{Potter2004}.
\end{prob}

Le point essentiel est que la partie
d'un ensemble de choix $C$ contenue
dans $\bigcup A$ est l'image d'une
fonction de choix pour $A$ ; inversement,
toute telle image est un ensemble de
choix pour $A$.\footnote{Note de l'édition
française : l'original identifie directement
$C$ à l'image d'une fonction de choix.
La définition permet pourtant à $C$
de contenir des objets hors de
$\bigcup A$ ; il faut prendre
$C\cap\bigcup A$.}
Pour justifier l'axiome du choix,
nous pouvons donc tenter de justifier
sa formulation équivalente en termes
d'existence d'ensembles de choix.
C'est ce que nous allons faire.

Supposons que les !!{element}s de $A$
soient disjoints et non vides.
D'après \stageshier{}
(voir la \olref[z][story]{sec}),
$A$ se forme à une certaine étape~$S$.
Tous les !!{element}s de $\bigcup A$
sont alors disponibles avant l'étape
$S$. Or, d'après \stagesacc{},
pour \emph{tous} les ensembles formés
avant~$S$, il se forme un ensemble
dont les éléments sont exactement
ceux-là. Autrement dit, toute
collection \emph{possible}
d'ensembles déjà disponibles existera
à l'étape~$S$. Il est certainement
\emph{possible} de sélectionner
des objets qui formeraient un ensemble
de choix pour~$A$ : il s'agirait
simplement d'un sous-ensemble bien
particulier de $\bigcup A$.
Un tel ensemble de choix existe
donc, comme souhaité.

Voilà une tentative \emph{très}
rapide de justifier l'axiome
du choix à partir de considérations
intrinsèques. Pour approfondir
cette idée, vous pouvez lire
l'élégant développement qu'en donne
Potter (\citeyear[\S14.8]{Potter2004}).

\end{document}
